\documentclass[12pt]{article}
\usepackage[utf8]{inputenc}
\usepackage[T1]{fontenc}
\usepackage{hyperref}
\title{Trieng Language Resources for K-12 Curriculum Development}
\author{Compiled by Claude Sonnet 5.5 (claude-sonnet-5-5) for Andr\'as Kornai}
\date{}
\begin{document}
\maketitle

\section{Introduction}
Trieng (ISO 639-3: stg; Glottocode trie1243), called Gie Trieng in Vietnamese, is a Bahnaric language of the Austroasiatic family. It is spoken mainly in Kon Tum and Quang Nam provinces, Vietnam. Wikipedia (in its article ``Tariang language'') gives about 45,000 speakers. It also notes that the ISO codes tdf (Talieng) and tgr (Tareng) cover related varieties, and these must not be confused with stg. Vietnamese sources give differing population counts for the ethnic group, such as 63,322 in the 2019 survey. The searches found no textbooks, dictionaries, primers or scripture for the language.

\section{Primary Educational Resources}
None were found in the search results.

\section{Linguistic and Reference Materials}
\begin{itemize}
\item Wikipedia, ``Tariang language'': \url{https://en.wikipedia.org/wiki/Tariang_language}
\item Numeral systems database, Trieng (Kon Tum): \url{https://lingweb.eva.mpg.de/channumerals/Trieng.htm}
\end{itemize}

\section{Religious and Community Resources}
\begin{itemize}
\item Joshua Project, Trieng in Vietnam (it says the group writes in Latin script and is largely Christian): \url{https://joshuaproject.net/people_groups/20454}
\item Nhan Dan, article on the Gie Trieng people (in Vietnamese): \url{https://nhandan.vn/dan-toc-gie-trieng-post723929.html}
\item Bao Van Hoa, ``Nguoi Gie Trieng gin giu mach nguon van hoa'' (in Vietnamese, on cultural preservation): \url{https://baovanhoa.vn/dan-toc-ton-giao/nguoi-gie-trieng-gin-giu-mach-nguon-van-hoa-139599.html}
\end{itemize}

\section{Digital and Audio Resources}
No audio, video or corpus resources were found.

\section{Curriculum Development Context}
The reference description says the language is not known to be taught in schools. Vietnamese is the language of instruction in the area. Any native-language curriculum would first need a vocabulary and an orthography, building on the Latin script already used. This should be done together with the community and with Vietnamese education authorities. A search in Vietnamese for educational materials in the language returned nothing useful.
\end{document}